\subsection{Original PL constructions retaining finite monomial ideal presentations}

\begin{theorem}[Actual original PL ideals with finite homogeneous parameter monomials]
\label{local:actual-pl-finite-monomial-models}
\lean{EnvelopeProperty.ProjSubvariety.IsPL.exists_monomial_vertical_productIdeal_models}
\leanok
\uses{local:actual-finite-parameter-monomial-powers}
For any original projective subvariety $X$ over any field $k$, and any
ORIGINAL PL function $\varphi$ on $X^{\rm an}$, there are $r>0$, an integer
$C$, and actual finitely generated homogeneous ideals $I,J$ in the
original polynomial projective model over $k[t]$ such that each contains
a positive power of $t$, each WHOLE underlying ideal is spanned by a
finite family of ORIGINAL parameter monomials with homogeneous base
forms, and the TRUE original model logarithms at every original Gauss
product point are finite and satisfy
$r\varphi=\log|I|-\log|J|+C$.
All finite monomial presentations are produced from the ORIGINAL PL data,
not supplied as an invariance or presentation hypothesis.
\end{theorem}
\begin{proof}
Use the ORIGINAL PL decomposition as a difference of two actual finite
FS functions. For each, the constructed genuine vertical
product ideal retains its actual finite homogeneous parameter-monomial
family, a positive parameter power, and its exact scaled finite model
logarithm. Cross-multiply the two positive integer scales, taking the
corresponding genuine ideal powers. The preceding finite monomial-power
construction retains actual finite homogeneous monomial families for
BOTH resulting WHOLE ideals. The parameter powers remain positive after
multiplying their exponents by the opposite positive scales. Finite
generation follows over the original Noetherian polynomial ring.
The true ideal-power logarithm formula multiplies each original finite
logarithm by the opposite scale. Their difference gives exactly the
original PL function at the product scale and integer shift
$r_2C_1-r_1C_2$. All asserted finiteness follows from these actual finite
real formulas, without an additional nonvanishing premise.
\end{proof}
